\documentstyle[% 


righttag% 


,11pt% 


,amssymb% 


,russian%               remove in Eng. 


,izvru%                 Set izven% in Eng. 


]{amsart} 


 


%  


\newcommand{\down}[2]{\underset{#1}{#2}{}} 


\newcommand{\const}{\operatorname{const}} 


\newcommand{\Aff}{\operatorname{Aff}} 


\newcommand{\tts}[1]{} 


 


%\hoffset=-15mm%                  For Eng. 


 


\setcounter{page}{3} 


\god{2000} 


\nomer{7~(458)} %                        Remove in Eng. 


%\nomer{Vol.\,44, No.\,7}%               Set in Eng. 


%\str{\pageref{begin}--\pageref{end}}%  Set in Eng. 


\udk{512.812\,:\,514.754} 


 


\author{\it А.В.\,Винник, С.Г.\,Лейко} 


% NB:   ^^ remove in Eng. 


\title{Изопериметрические экстремали поворота на двумерных 


связных группах Ли с инвариантными римановыми метриками} 


 


\date{} 


% \pagestyle{myheadings}%         Set in Eng. 


\pagestyle{plain} %              Remove in Eng. 


\begin{document} 


% \thispagestyle{plain}%          Set in Eng. 


\maketitle 


\label{begin} 


 


На пространстве параметризованных кривых в двумерном римановом 


многообразии $(M,g)$ определены два естественных функционала: 


функционал 


длины 


$$ 


l[\gamma]=\int^{t_1}_{t_0}\sqrt{g(\Dot\gamma,\Dot\gamma)}\,dt, 


$$ 


где $\Dot\gamma$ --- касательный вектор параметризованной кривой 


$\gamma$, и 


функционал абсолютного поворота 


$$ 


\Theta[\gamma]=\int^{l_1}_{l_0}k_1(l)d l, 


$$ 


где $l$ --- длина дуги на кривой $\gamma$, $k_1(l)$


--- первая кривизна кривой. 


 


В [1], [2] было показано, что стационарные кривые изопериметрической 


вариационной задачи с фиксированными концами 


$$ 


\text{extremum }\Theta[\gamma],\ \; \gamma(t_0)=p_0, 


\ \; \gamma(t_1)=p_1,\ \; l[\gamma]=\const=\wh l, 


$$ 


либо удовлетворяют уравнению 


\begin{equation} 


K=\wh c k, 


\end{equation} 


где $K$ --- кривизна метрики $g$, $\wh c$ --- 


изопериметрическая постоянная, зависящая от $\wh l$, либо являются 


геодезическими. 


 


Кривые, удовлетворяющие уравнению (1), а также геодезические кривые 


названы 


изопериметрическими экстремалями поворота (ИЭП) [1]. 


Геодезические естественно назвать 


тривиальными ИЭП. 


 


Для многообразия нулевой кривизны $K=0$ всякая допустимая кривая 


является 


при $\wh c=0$ его ИЭП. Если $K\ne0$, то (1) можно представить в виде 


\begin{equation} 


k=c K,\tag{$1'$} 


\end{equation} 


где $k$ вычисляется по формуле 


$$ 


k=\frac{(g_{11}g_{22}-g^2_{12})^{1/2}} 


{(g_{i j}\Dot u^i\Dot u^j)^{3/2}}[\Dot u^1(\Ddot u^2+ 


\Gamma^2_{i j}\Dot u^i\Dot u^j)- 


\Dot u^2(\Ddot u^1+\Gamma^1_{i j} 


\Dot u^i\Dot u^j)], 


$$ 


и случаем $c=0$ охватываются также тривиальные ИЭП. 


Тем самым изучение ИЭП содержательно при $K\ne0$. 


 


Известно [3], что если $M^2$ --- связная двумерная группа Ли, 


то она изоморфна одной из следующих групп: 


\begin{itemize} 


\item[I] (абелевы группы) 1) $\Bbb R^2$, 2) $T^1\times\Bbb R$, 3) $T^2$; 


\item[II] (не абелева группа) 4) $\Aff^0\Bbb R^1$ --- группа собственных 


аффинных 


преобразований прямой $x^1=b^2x+b^1$, $b^2>0$. 


\end{itemize} 


\noindent{}Здесь $T^1=\{z\in C\mid z=e^{i\varphi}\}$, $T^2=T^1\times T^1$


--- соответственно одномерный и двумерный торы. 


 


В случае I групповые функции имеют вид $f^1(a,b)=a^1+b^1$, 


$f^2(a,b)=a^2+b^2$ 


и $e=(0,0)$ --- единица группы. Отсюда получаем


координатное представление дифференциала левого сдвига $L_a$:


$$ 


\down{*}{L}_a:L^i_j(a)= 


\frac{\partial f^i(a,b)}{\partial b^j}\bigg|_{b=e}=\delta^i_j. 


$$ 


Компоненты левоинвариантной метрики ($\down{*}{L}g=g$) вычисляются по 


формуле [4]


\begin{equation} 


g_{i j}(a)=g_{p q}(e)V^p_i(a)V^q_j(a),\quad 


V^i_j L^j_k=\delta^i_k, 


\end{equation} 


где $g_{p q}(e)$ (компоненты метрики в единице группы) могут выбираться 


произвольно. Возьмем их канонические значения $g_{p q}(e)=\delta_{p q}$ 


и в дальнейшем 


соответствующие римановы инвариантные метрики $g$ на $M^2$ 


будем называть каноническими. Так как для случая I 


$V^i_j(a)=\delta^i_j$, то каноническая левоинвариантная риманова 


метрика имеет компоненты 


$g_{i j}=\delta_{i j}$. 


 


Аналогично рассматриваются правые сдвиги $R_b:x\to x b$, и получается 


координатное 


представление дифференциала 


$$ 


\down{*}{R}_b:R^i_j(b)= 


\frac{\partial f^i(a,b)}{\partial a^j}\bigg|_{a=e}=\delta^i_j. 


$$ 


Компоненты правоинвариантной метрики ($\down{*}{R}g=g$) вычисляются по 


формуле 


\begin{equation} 


g_{i j}(a)=g_{p q}(e)U^p_i(a)U^q_j(a),\quad 


U^i_j L^j_k=\delta^i_k. 


\end{equation} 


В случае I\; $U^i_j(b)=\delta^i_j$ и для канонической 


правоинвариантной 


метрики получаем $g_{i j}(a)=\delta_{i j}$. 


Таким образом, существует биинвариантная каноническая метрика 


$g_{i j}(a)=\delta_{i j}$ и 


$(M^2,g)$ с этой метрикой является римановым многообразием нулевой 


кривизны $K=0$. Отметим, что $M$ может быть гомеоморфно плоскости, 


цилиндру или тору (~1), 2), 3) соответственно). 


 


В случае II групповые функции имеют вид 


$$ 


f^1(a,b)=a^1+a^2b^1,\quad 


f^2(a,b)=a^2b^2,\quad a^2,b^2>0. 


$$ 


Отсюда $e=(0,1)$, и 


\begin{alignat*}{2} 


L^i_j(a)&=\begin{pmatrix} 


a^2 & 0 \\ 0 & a^2 


\end{pmatrix} 


=a^2\delta^h_i, &\quad V^i_j(a)&=\frac{1}{a^2}\delta^i_j,\\ 


R^i_j(b)&=\begin{pmatrix} 


1 & b^1 \\ 0 & b^2 


\end{pmatrix} 


,& \quad U^i_j(b)&= 


\begin{pmatrix} 


1 & -\dfrac{b^1}{b^2}\\[3mm] 0 & \dfrac{1}{b^2} 


\end{pmatrix} 


, 


\end{alignat*} 


поэтому левоинвариантная каноническая метрика в 


соответствии с (2) имеет компоненты 


$$ 


g_{i j}(a)=\delta_{p q}\frac{1}{a^2} 


\delta^p_i\frac{1}{a^2}\delta^q_j= 


\frac{1}{(a^2)^2}\delta_{i j}. 


$$ 


Полученная метрика имеет постоянную кривизну $K=-1$, и риманово 


многообразие $(M,g)$ может быть отождествлено с моделью Пуанкаре 


плоскости Лобачевского на верхней полуплоскости. 


Из $(1')$ вытекает, что в данном случае ИЭП --- геодезические 


окружности 


плоскости Лобачевского. Известно [5], что 


последние в модели Пуанкаре состоят из двух семейств: 


\begin{itemize} 


\item[1.] прямые $y=k x+b$, $k_1=\frac{1}{\sqrt{1+k^2}}$ или $x=x_0$, 


$k_1=0$; 


\item[2.] окружности (или их части) $(x-x_0)^2+(y-y_0)^2=r^2$, 


$k_1=y_0/r$. 


\end{itemize} 


\noindent{}В силу (3) в данном случае правоинвариантная метрика имеет 


компоненты 


$$ 


(g_{i j}(b))= 


\begin{pmatrix} 


1 & - \dfrac{b^1}{b^2}\\[3mm] 


-\dfrac{b^1}{b^2} & \dfrac{1+(b^1)^2}{(b^2)^2} 


\end{pmatrix} 


. 


$$ 


Эта метрика имеет также постоянную отрицательную кривизну 


$K=-1$.


Как видно из вышеизложенного, на группе $\Aff^0\Re^1$ не существует 


биинвариантной канонической римановой метрики. 


 


Отметим в заключение, что выбор не канонического 


значения $g_{i j}(e)$ приводит к инвариантным римановым метрикам, 


которые будут получаться из канонических инвариантных метрик 


аффинным преобразованием координат. Тем самым получено 


описание ИЭП с точностью до 


таких преобразований. 


 


\begin{thebibliography}{9} 


 


\bibitem{1} 


Лейко С.Г. {\em Вариационные задачи для функционалов поворота и 


спин-отображения псевдоримановых пространств} // 


Изв. вузов. Математика. -- 1990. -- \No\,10. -- C.\,9--17. 


\bibitem{2} 


Лейко С.Г. {\em Изопериметрические экстремали поворота на 


поверхностях в евклидовом пространстве $E^3$} // 


Изв. вузов. Математика. -- 1996. -- \No\,6. -- C.\,25--31. 


\bibitem{3} 


Винберг Э.Б., Горбацевич В.В., Онищик А.Л. {\em Строение групп и алгебр 


Ли} // Итоги науки и техн. ВИНИТИ. Современ. пробл. матем. 


-- 1990. -- Т.\,41. -- 254 с. 


\bibitem{4} 


Шапуков Б.Н. {\em Задачи по группам Ли}. Учеб. пособие. 


-- Изд-во. Казанск. ун-та, 1989. -- 152~с. 


\bibitem{5} 


Букреев Б.Я. {\em Планиметрия Лобачевского в аналитическом 


изложении}. -- М.-Л.: ГИТЛ, 1951. -- 127~с. 


 


\end{thebibliography} 


 


\vskip 2cm 


 


\post{\it Одесский государственный}{\it Поступила} 


\post{\it университет $($Украина$)$}{08.06.1998} 


 


\label{end} 


\end{document} 


 


